Praca dyplomowa dotyczy zaprojektowania i weryfikacji systemu sterowania dla bezzałogowego pojazdu podwodnego. W początkowej fazie system przeznaczony jest dla pojazdu zdalnie sterowanego (ROV), jednak planowana jest jego adaptacja do wersji autonomicznej (AUV). Głównym wyzwaniem inżynierskim podjętym w pracy jest analityczne wyprowadzenie nieliniowego modelu matematycznego dynamiki statku w sześciu stopniach swobody, uwzględniającego specyfikę środowiska wodnego, w tym zjawiska mas dodanych oraz tłumienia hydrodynamicznego. W ramach pracy zastosowane zsotaną klasyczne struktury regulacji (PID) oraz algorytmy uwzględniające indywidualną dynamikę obiektu. Z uwagi na silne tłumienie fal elektromagnetycznych pod wodą i brak sygnału GPS, pętlę sprzężenia zwrotnego oparta będzie na danych z czujników inercyjnych (IMU) oraz czujnika ciśnienia. Oprogramowanie systemu sterowania zostanie opracowane w oparciu o śwodowisko ROS 2 i zweryfikowane z wykorzystaniem symulatora Gazebo. Całość architektury programowej zostanie przygotowana do bezpośredniego wdrożenia na fizycznym sprzęcie, np. Raspberry Pi.